\section{Introduction}

Egocentric video understanding has emerged as a central challenge in embodied AI, with applications
ranging from assistive technology to augmented reality and robot learning from human
demonstration~\cite{grauman2022ego4d,grauman2024egoexo4d,wang2023holoassist}.
Unlike third-person video, egocentric footage captures the world from the perspective of an active
agent, embedding rich signals about intent, attention, and physical interaction directly within the
visual stream.
A key indicator of this intentionality is \emph{gaze}—where the camera wearer directs their
attention~\cite{li2018egtea,huang2018predicting}—combined with the motion of their hands as they
manipulate objects in the scene~\cite{nagarajan2020ego}.

Recent vision-language models (VLMs)~\cite{wang2024qwen2vl,lin2023videollava} have demonstrated
strong video question-answering (VideoQA) capabilities, but their applicability to egocentric video
is constrained by a fundamental bottleneck: the number of visual tokens that can be processed
within the context window.
A 128-frame egocentric clip processed at full spatial resolution produces tens of thousands of
visual patch tokens, far exceeding the token budgets of modern VLMs.
Existing approaches address this by uniformly subsampling frames or spatially pooling features,
discarding information without regard for where the agent's attention was directed.
Prior work on efficient vision transformers has shown that not all tokens contribute equally to
downstream predictions~\cite{rao2021dynamicvit,liang2022evit}, and that merging redundant tokens
can substantially reduce computation without sacrificing
accuracy~\cite{bolya2023tome}.
However, these methods operate on general visual salience and are agnostic to the rich behavioral
signals available in egocentric settings.

We argue that gaze and hand trajectory data—routinely recorded by egocentric devices—provide a
principled signal for \emph{which} visual tokens carry the most task-relevant information.
Moments of sustained fixation and bimanual manipulation correlate strongly with the objects,
actions, and spatial regions that egocentric QA tasks ask about~\cite{li2018egtea,mangalam2024egoschema}.
Yet to our knowledge, no prior work has exploited this signal to guide visual token selection
inside a VLM's attention layers.

In this work, we propose \textbf{TrajGazeMerge}, a framework that uses egocentric gaze and hand
trajectories to identify and retain the most informative visual tokens, discarding or folding the
remainder before VLM processing.
Our approach consists of three components:

\begin{enumerate}[leftmargin=*,itemsep=2pt,topsep=2pt]
\item \textbf{TrajGaze Encoder.} A query-conditioned spatio-temporal network that ingests
    per-frame gaze positions, left- and right-hand trajectories, and their pairwise interaction
    features, then cross-attends these trajectory signals to the spatial patch grid of a frozen
    visual backbone~\cite{oquab2023dinov2} to produce a per-patch importance score.

\item \textbf{Bipartite Token Merging.} Given per-patch scores, we designate the top-scoring
    patches as \emph{receivers} and the remaining patches as \emph{sources}.
    Each source is merged into its cosine-nearest receiver via a score-weighted average~\cite{bolya2023tome},
    preserving information from discarded regions rather than simply dropping them.
    We retain only 10\% of the original token count, reducing VLM context length by $10\times$.

\item \textbf{Joint LoRA Fine-tuning.} We fine-tune the TrajGaze encoder and lightweight
    LoRA~\cite{hu2022lora} adapters on the VLM jointly, using a combination of cross-entropy
    loss on the ground-truth answer and a KL-divergence regularization toward a teacher VLM
    operating on the full token set.
\end{enumerate}

We evaluate on \textbf{StreamGaze\_v2}, an egocentric VideoQA benchmark built on
EGTEA~\cite{li2018egtea}, EgoExo4D~\cite{grauman2024egoexo4d}, and
HoloAssist~\cite{wang2023holoassist}, spanning eight task categories that probe spatial,
temporal, and gaze-predictive reasoning.
TrajGazeMerge achieves \textbf{64.45\%} accuracy on the EGTEA test split, outperforming a
random-10\% token selection baseline (61.98\%) and the full-token zero-shot VLM baseline.
Per-task analysis reveals that trajectory-guided selection provides the largest gains on tasks
requiring temporal gaze reasoning (\emph{past\_gaze\_sequence\_matching}: +7.8~pp) and
non-fixated object identification (\emph{past\_non\_fixated\_object\_identification}: +7.4~pp).

Our work demonstrates that egocentric behavioral signals—gaze and hand motion—are not merely
outputs to be predicted, but \emph{inputs} that can actively guide efficient visual processing
in VLMs, opening a new direction for resource-efficient egocentric video understanding.
